% paper101-judgment-fiber-bundles.tex --- Judgment Fiber Bundles: Trust as a Geometric Connection in Multi-Agent AI Verification
%   Paper 101 of the Judgment Geometry series.
% Compile: pdflatex paper101-judgment-fiber-bundles
\documentclass[11pt]{article}
\usepackage{lmodern}
% jugeo-common.tex — shared preamble for all Judgment Geometry papers
% Include via: % jugeo-common.tex — shared preamble for all Judgment Geometry papers
% Include via: % jugeo-common.tex — shared preamble for all Judgment Geometry papers
% Include via: \input{jugeo-common}

% ── Packages ────────────────────────────────────────────────────────────
\usepackage{amsmath,amssymb,amsthm,mathtools}
\usepackage{tikz-cd}
\usepackage{listings}
\usepackage{xcolor}
\usepackage{xspace}
\usepackage{booktabs}
\usepackage{multirow}
\usepackage{enumitem}
\usepackage[expansion=false]{microtype}
\usepackage{hyperref}
\usepackage{cleveref}
% \usepackage{thmtools}  % disabled: not available in this TeX installation
\usepackage{float}
\usepackage{caption}
\usepackage{subcaption}
\usepackage{tabularx}
\usepackage{stmaryrd}       % \llbracket, \rrbracket
\usepackage{bbm}            % \mathbbm{1}

% ── Page geometry ───────────────────────────────────────────────────────
\usepackage[margin=1in]{geometry}

% ── Theorem environments ────────────────────────────────────────────────
\theoremstyle{plain}
\newtheorem{theorem}{Theorem}[section]
\newtheorem{lemma}[theorem]{Lemma}
\newtheorem{proposition}[theorem]{Proposition}
\newtheorem{corollary}[theorem]{Corollary}

\theoremstyle{definition}
\newtheorem{definition}[theorem]{Definition}
\newtheorem{example}[theorem]{Example}
\newtheorem{construction}[theorem]{Construction}

\theoremstyle{remark}
\newtheorem{remark}[theorem]{Remark}
\newtheorem{notation}[theorem]{Notation}

% ── Python listings style ──────────────────────────────────────────────
\definecolor{jugeoBlue}{HTML}{2563EB}
\definecolor{jugeoGreen}{HTML}{059669}
\definecolor{jugeoRed}{HTML}{DC2626}
\definecolor{jugeoPurple}{HTML}{7C3AED}
\definecolor{jugeoGray}{HTML}{6B7280}
\definecolor{jugeoBg}{HTML}{F8FAFC}

\lstdefinestyle{jugeo-python}{
  language=Python,
  basicstyle=\ttfamily\small,
  keywordstyle=\color{jugeoBlue}\bfseries,
  stringstyle=\color{jugeoGreen},
  commentstyle=\color{jugeoGray}\itshape,
  numberstyle=\tiny\color{jugeoGray},
  backgroundcolor=\color{jugeoBg},
  numbers=left,
  numbersep=6pt,
  frame=single,
  framerule=0.4pt,
  rulecolor=\color{jugeoGray!40},
  breaklines=true,
  breakatwhitespace=true,
  showstringspaces=false,
  tabsize=4,
  xleftmargin=1.5em,
  framexleftmargin=1.5em,
  morekeywords={dataclass,frozen,field,Enum,IntEnum,auto,Any,
                Mapping,Sequence,Optional,match,case},
  literate={→}{{$\to$}}1 {←}{{$\leftarrow$}}1
           {≤}{{$\leq$}}1 {≥}{{$\geq$}}1 {≠}{{$\neq$}}1
           {∈}{{$\in$}}1 {∀}{{$\forall$}}1 {∃}{{$\exists$}}1
           {⊕}{{$\oplus$}}1 {⊖}{{$\ominus$}}1,
  escapeinside={(*@}{@*)},
}
\lstset{style=jugeo-python}

\lstdefinestyle{jugeo-output}{
  basicstyle=\ttfamily\small,
  backgroundcolor=\color{jugeoBg},
  frame=single,
  framerule=0.4pt,
  rulecolor=\color{jugeoGray!40},
  breaklines=true,
  showstringspaces=false,
  xleftmargin=1.5em,
  framexleftmargin=0.5em,
  numbers=none,
}

% ── JuGeo-specific macros ──────────────────────────────────────────────

% Judgment 8-tuple
\newcommand{\J}{\mathcal{J}}
\newcommand{\Jud}[8]{(#1,\,#2,\,#3,\,#4,\,#5,\,#6,\,#7,\,#8)}
\newcommand{\JudShort}[1]{\J_{#1}}

% Components
\newcommand{\coord}{\mathsf{c}}            % coordinate
\newcommand{\prop}{\varphi}                 % proposition
\newcommand{\carrier}{\mathcal{A}}          % carrier type
\newcommand{\evid}{\mathcal{E}}             % evidence bundle
\newcommand{\oblig}{\mathcal{O}}            % obligations
\newcommand{\obstr}{\mathcal{B}}            % obstructions
\newcommand{\trust}{\mathcal{T}}            % trust annotation
\newcommand{\prov}{\Pi}                     % provenance

% Trust algebra
\newcommand{\Talg}{\mathfrak{T}}            % trust ordered algebra
\newcommand{\Eadm}{\mathcal{E}_{\mathrm{adm}}}
\newcommand{\tleq}{\preceq}                 % trust ordering
\newcommand{\tjoin}{\oplus}                 % conservative join
\newcommand{\tatten}{\ominus}               % attenuation
\newcommand{\tpromote}{\uparrow_\pi}        % promotion
\newcommand{\tdemote}{\downarrow_\chi}       % demotion

% Trust tiers
\newcommand{\Tcontra}{\ensuremath{\mathsf{CONTRADICTED}}}
\newcommand{\Tunver}{\ensuremath{\mathsf{UNVERIFIED}}}
\newcommand{\Tcopilot}{\ensuremath{\mathsf{COPILOT}}}
\newcommand{\Toracle}{\ensuremath{\mathsf{ORACLE}}}
\newcommand{\Truntime}{\ensuremath{\mathsf{RUNTIME}}}
\newcommand{\Tsolver}{\ensuremath{\mathsf{SOLVER}}}
\newcommand{\Tproof}{\ensuremath{\mathsf{PROOF}}}

% Site and geometry
\newcommand{\Site}{\mathcal{S}}             % semantic site
\newcommand{\Cov}{\mathrm{Cov}}            % covering family
\newcommand{\Ob}{\mathrm{Ob}}              % objects
\newcommand{\Mor}{\mathrm{Mor}}            % morphisms
\newcommand{\res}{\mathrm{res}}            % restriction
\newcommand{\Cech}{\check{C}}              % Čech complex
\newcommand{\Hcoh}[1]{H^{#1}}             % cohomology group

% Descent
\newcommand{\descent}{\mathrm{desc}}
\newcommand{\glue}{\mathrm{glue}}
\newcommand{\overlap}[2]{#1 \cap #2}

% Semantic moves
\newcommand{\VERIFY}{\textsc{Verify}}
\newcommand{\CONSTRUCT}{\textsc{Construct}}
\newcommand{\NORMALIZE}{\textsc{Normalize}}
\newcommand{\GLUE}{\textsc{Glue}}
\newcommand{\RESTRICT}{\textsc{Restrict}}
\newcommand{\TRANSPORT}{\textsc{Transport}}
\newcommand{\REPAIR}{\textsc{Repair}}
\newcommand{\PROMOTE}{\textsc{Promote}}
\newcommand{\CHALLENGE}{\textsc{Challenge}}
\newcommand{\ESCALATE}{\textsc{Escalate}}

% Fragments
\newcommand{\QFLIA}{\texttt{QF\_LIA}}
\newcommand{\QFLRA}{\texttt{QF\_LRA}}
\newcommand{\QFBV}{\texttt{QF\_BV}}
\newcommand{\QFUF}{\texttt{QF\_UF}}

% Misc
\newcommand{\Python}{\textsc{Python}}
\newcommand{\JuGeo}{\textsc{JuGeo}}
\newcommand{\LEAN}{\textsc{Lean}\,4}
\newcommand{\Fstar}{F$^{\star}$}
\newcommand{\Coq}{\textsc{Coq}}
\newcommand{\Dafny}{\textsc{Dafny}}

% Common shortcuts
\newcommand{\ie}{i.e.\@\xspace}
\newcommand{\eg}{e.g.\@\xspace}
\newcommand{\cf}{cf.\@\xspace}
\newcommand{\wrt}{w.r.t.\@\xspace}

% Evaluation shortcuts
\newcommand{\numcases}{300}
\newcommand{\numspec}{100}
\newcommand{\numeq}{100}
\newcommand{\numbug}{100}
\newcommand{\medtime}{6.5\,ms}
\newcommand{\totaltime}{1.949\,s}


% ── Packages ────────────────────────────────────────────────────────────
\usepackage{amsmath,amssymb,amsthm,mathtools}
\usepackage{tikz-cd}
\usepackage{listings}
\usepackage{xcolor}
\usepackage{xspace}
\usepackage{booktabs}
\usepackage{multirow}
\usepackage{enumitem}
\usepackage[expansion=false]{microtype}
\usepackage{hyperref}
\usepackage{cleveref}
% \usepackage{thmtools}  % disabled: not available in this TeX installation
\usepackage{float}
\usepackage{caption}
\usepackage{subcaption}
\usepackage{tabularx}
\usepackage{stmaryrd}       % \llbracket, \rrbracket
\usepackage{bbm}            % \mathbbm{1}

% ── Page geometry ───────────────────────────────────────────────────────
\usepackage[margin=1in]{geometry}

% ── Theorem environments ────────────────────────────────────────────────
\theoremstyle{plain}
\newtheorem{theorem}{Theorem}[section]
\newtheorem{lemma}[theorem]{Lemma}
\newtheorem{proposition}[theorem]{Proposition}
\newtheorem{corollary}[theorem]{Corollary}

\theoremstyle{definition}
\newtheorem{definition}[theorem]{Definition}
\newtheorem{example}[theorem]{Example}
\newtheorem{construction}[theorem]{Construction}

\theoremstyle{remark}
\newtheorem{remark}[theorem]{Remark}
\newtheorem{notation}[theorem]{Notation}

% ── Python listings style ──────────────────────────────────────────────
\definecolor{jugeoBlue}{HTML}{2563EB}
\definecolor{jugeoGreen}{HTML}{059669}
\definecolor{jugeoRed}{HTML}{DC2626}
\definecolor{jugeoPurple}{HTML}{7C3AED}
\definecolor{jugeoGray}{HTML}{6B7280}
\definecolor{jugeoBg}{HTML}{F8FAFC}

\lstdefinestyle{jugeo-python}{
  language=Python,
  basicstyle=\ttfamily\small,
  keywordstyle=\color{jugeoBlue}\bfseries,
  stringstyle=\color{jugeoGreen},
  commentstyle=\color{jugeoGray}\itshape,
  numberstyle=\tiny\color{jugeoGray},
  backgroundcolor=\color{jugeoBg},
  numbers=left,
  numbersep=6pt,
  frame=single,
  framerule=0.4pt,
  rulecolor=\color{jugeoGray!40},
  breaklines=true,
  breakatwhitespace=true,
  showstringspaces=false,
  tabsize=4,
  xleftmargin=1.5em,
  framexleftmargin=1.5em,
  morekeywords={dataclass,frozen,field,Enum,IntEnum,auto,Any,
                Mapping,Sequence,Optional,match,case},
  literate={→}{{$\to$}}1 {←}{{$\leftarrow$}}1
           {≤}{{$\leq$}}1 {≥}{{$\geq$}}1 {≠}{{$\neq$}}1
           {∈}{{$\in$}}1 {∀}{{$\forall$}}1 {∃}{{$\exists$}}1
           {⊕}{{$\oplus$}}1 {⊖}{{$\ominus$}}1,
  escapeinside={(*@}{@*)},
}
\lstset{style=jugeo-python}

\lstdefinestyle{jugeo-output}{
  basicstyle=\ttfamily\small,
  backgroundcolor=\color{jugeoBg},
  frame=single,
  framerule=0.4pt,
  rulecolor=\color{jugeoGray!40},
  breaklines=true,
  showstringspaces=false,
  xleftmargin=1.5em,
  framexleftmargin=0.5em,
  numbers=none,
}

% ── JuGeo-specific macros ──────────────────────────────────────────────

% Judgment 8-tuple
\newcommand{\J}{\mathcal{J}}
\newcommand{\Jud}[8]{(#1,\,#2,\,#3,\,#4,\,#5,\,#6,\,#7,\,#8)}
\newcommand{\JudShort}[1]{\J_{#1}}

% Components
\newcommand{\coord}{\mathsf{c}}            % coordinate
\newcommand{\prop}{\varphi}                 % proposition
\newcommand{\carrier}{\mathcal{A}}          % carrier type
\newcommand{\evid}{\mathcal{E}}             % evidence bundle
\newcommand{\oblig}{\mathcal{O}}            % obligations
\newcommand{\obstr}{\mathcal{B}}            % obstructions
\newcommand{\trust}{\mathcal{T}}            % trust annotation
\newcommand{\prov}{\Pi}                     % provenance

% Trust algebra
\newcommand{\Talg}{\mathfrak{T}}            % trust ordered algebra
\newcommand{\Eadm}{\mathcal{E}_{\mathrm{adm}}}
\newcommand{\tleq}{\preceq}                 % trust ordering
\newcommand{\tjoin}{\oplus}                 % conservative join
\newcommand{\tatten}{\ominus}               % attenuation
\newcommand{\tpromote}{\uparrow_\pi}        % promotion
\newcommand{\tdemote}{\downarrow_\chi}       % demotion

% Trust tiers
\newcommand{\Tcontra}{\ensuremath{\mathsf{CONTRADICTED}}}
\newcommand{\Tunver}{\ensuremath{\mathsf{UNVERIFIED}}}
\newcommand{\Tcopilot}{\ensuremath{\mathsf{COPILOT}}}
\newcommand{\Toracle}{\ensuremath{\mathsf{ORACLE}}}
\newcommand{\Truntime}{\ensuremath{\mathsf{RUNTIME}}}
\newcommand{\Tsolver}{\ensuremath{\mathsf{SOLVER}}}
\newcommand{\Tproof}{\ensuremath{\mathsf{PROOF}}}

% Site and geometry
\newcommand{\Site}{\mathcal{S}}             % semantic site
\newcommand{\Cov}{\mathrm{Cov}}            % covering family
\newcommand{\Ob}{\mathrm{Ob}}              % objects
\newcommand{\Mor}{\mathrm{Mor}}            % morphisms
\newcommand{\res}{\mathrm{res}}            % restriction
\newcommand{\Cech}{\check{C}}              % Čech complex
\newcommand{\Hcoh}[1]{H^{#1}}             % cohomology group

% Descent
\newcommand{\descent}{\mathrm{desc}}
\newcommand{\glue}{\mathrm{glue}}
\newcommand{\overlap}[2]{#1 \cap #2}

% Semantic moves
\newcommand{\VERIFY}{\textsc{Verify}}
\newcommand{\CONSTRUCT}{\textsc{Construct}}
\newcommand{\NORMALIZE}{\textsc{Normalize}}
\newcommand{\GLUE}{\textsc{Glue}}
\newcommand{\RESTRICT}{\textsc{Restrict}}
\newcommand{\TRANSPORT}{\textsc{Transport}}
\newcommand{\REPAIR}{\textsc{Repair}}
\newcommand{\PROMOTE}{\textsc{Promote}}
\newcommand{\CHALLENGE}{\textsc{Challenge}}
\newcommand{\ESCALATE}{\textsc{Escalate}}

% Fragments
\newcommand{\QFLIA}{\texttt{QF\_LIA}}
\newcommand{\QFLRA}{\texttt{QF\_LRA}}
\newcommand{\QFBV}{\texttt{QF\_BV}}
\newcommand{\QFUF}{\texttt{QF\_UF}}

% Misc
\newcommand{\Python}{\textsc{Python}}
\newcommand{\JuGeo}{\textsc{JuGeo}}
\newcommand{\LEAN}{\textsc{Lean}\,4}
\newcommand{\Fstar}{F$^{\star}$}
\newcommand{\Coq}{\textsc{Coq}}
\newcommand{\Dafny}{\textsc{Dafny}}

% Common shortcuts
\newcommand{\ie}{i.e.\@\xspace}
\newcommand{\eg}{e.g.\@\xspace}
\newcommand{\cf}{cf.\@\xspace}
\newcommand{\wrt}{w.r.t.\@\xspace}

% Evaluation shortcuts
\newcommand{\numcases}{300}
\newcommand{\numspec}{100}
\newcommand{\numeq}{100}
\newcommand{\numbug}{100}
\newcommand{\medtime}{6.5\,ms}
\newcommand{\totaltime}{1.949\,s}


% ── Packages ────────────────────────────────────────────────────────────
\usepackage{amsmath,amssymb,amsthm,mathtools}
\usepackage{tikz-cd}
\usepackage{listings}
\usepackage{xcolor}
\usepackage{xspace}
\usepackage{booktabs}
\usepackage{multirow}
\usepackage{enumitem}
\usepackage[expansion=false]{microtype}
\usepackage{hyperref}
\usepackage{cleveref}
% \usepackage{thmtools}  % disabled: not available in this TeX installation
\usepackage{float}
\usepackage{caption}
\usepackage{subcaption}
\usepackage{tabularx}
\usepackage{stmaryrd}       % \llbracket, \rrbracket
\usepackage{bbm}            % \mathbbm{1}

% ── Page geometry ───────────────────────────────────────────────────────
\usepackage[margin=1in]{geometry}

% ── Theorem environments ────────────────────────────────────────────────
\theoremstyle{plain}
\newtheorem{theorem}{Theorem}[section]
\newtheorem{lemma}[theorem]{Lemma}
\newtheorem{proposition}[theorem]{Proposition}
\newtheorem{corollary}[theorem]{Corollary}

\theoremstyle{definition}
\newtheorem{definition}[theorem]{Definition}
\newtheorem{example}[theorem]{Example}
\newtheorem{construction}[theorem]{Construction}

\theoremstyle{remark}
\newtheorem{remark}[theorem]{Remark}
\newtheorem{notation}[theorem]{Notation}

% ── Python listings style ──────────────────────────────────────────────
\definecolor{jugeoBlue}{HTML}{2563EB}
\definecolor{jugeoGreen}{HTML}{059669}
\definecolor{jugeoRed}{HTML}{DC2626}
\definecolor{jugeoPurple}{HTML}{7C3AED}
\definecolor{jugeoGray}{HTML}{6B7280}
\definecolor{jugeoBg}{HTML}{F8FAFC}

\lstdefinestyle{jugeo-python}{
  language=Python,
  basicstyle=\ttfamily\small,
  keywordstyle=\color{jugeoBlue}\bfseries,
  stringstyle=\color{jugeoGreen},
  commentstyle=\color{jugeoGray}\itshape,
  numberstyle=\tiny\color{jugeoGray},
  backgroundcolor=\color{jugeoBg},
  numbers=left,
  numbersep=6pt,
  frame=single,
  framerule=0.4pt,
  rulecolor=\color{jugeoGray!40},
  breaklines=true,
  breakatwhitespace=true,
  showstringspaces=false,
  tabsize=4,
  xleftmargin=1.5em,
  framexleftmargin=1.5em,
  morekeywords={dataclass,frozen,field,Enum,IntEnum,auto,Any,
                Mapping,Sequence,Optional,match,case},
  literate={→}{{$\to$}}1 {←}{{$\leftarrow$}}1
           {≤}{{$\leq$}}1 {≥}{{$\geq$}}1 {≠}{{$\neq$}}1
           {∈}{{$\in$}}1 {∀}{{$\forall$}}1 {∃}{{$\exists$}}1
           {⊕}{{$\oplus$}}1 {⊖}{{$\ominus$}}1,
  escapeinside={(*@}{@*)},
}
\lstset{style=jugeo-python}

\lstdefinestyle{jugeo-output}{
  basicstyle=\ttfamily\small,
  backgroundcolor=\color{jugeoBg},
  frame=single,
  framerule=0.4pt,
  rulecolor=\color{jugeoGray!40},
  breaklines=true,
  showstringspaces=false,
  xleftmargin=1.5em,
  framexleftmargin=0.5em,
  numbers=none,
}

% ── JuGeo-specific macros ──────────────────────────────────────────────

% Judgment 8-tuple
\newcommand{\J}{\mathcal{J}}
\newcommand{\Jud}[8]{(#1,\,#2,\,#3,\,#4,\,#5,\,#6,\,#7,\,#8)}
\newcommand{\JudShort}[1]{\J_{#1}}

% Components
\newcommand{\coord}{\mathsf{c}}            % coordinate
\newcommand{\prop}{\varphi}                 % proposition
\newcommand{\carrier}{\mathcal{A}}          % carrier type
\newcommand{\evid}{\mathcal{E}}             % evidence bundle
\newcommand{\oblig}{\mathcal{O}}            % obligations
\newcommand{\obstr}{\mathcal{B}}            % obstructions
\newcommand{\trust}{\mathcal{T}}            % trust annotation
\newcommand{\prov}{\Pi}                     % provenance

% Trust algebra
\newcommand{\Talg}{\mathfrak{T}}            % trust ordered algebra
\newcommand{\Eadm}{\mathcal{E}_{\mathrm{adm}}}
\newcommand{\tleq}{\preceq}                 % trust ordering
\newcommand{\tjoin}{\oplus}                 % conservative join
\newcommand{\tatten}{\ominus}               % attenuation
\newcommand{\tpromote}{\uparrow_\pi}        % promotion
\newcommand{\tdemote}{\downarrow_\chi}       % demotion

% Trust tiers
\newcommand{\Tcontra}{\ensuremath{\mathsf{CONTRADICTED}}}
\newcommand{\Tunver}{\ensuremath{\mathsf{UNVERIFIED}}}
\newcommand{\Tcopilot}{\ensuremath{\mathsf{COPILOT}}}
\newcommand{\Toracle}{\ensuremath{\mathsf{ORACLE}}}
\newcommand{\Truntime}{\ensuremath{\mathsf{RUNTIME}}}
\newcommand{\Tsolver}{\ensuremath{\mathsf{SOLVER}}}
\newcommand{\Tproof}{\ensuremath{\mathsf{PROOF}}}

% Site and geometry
\newcommand{\Site}{\mathcal{S}}             % semantic site
\newcommand{\Cov}{\mathrm{Cov}}            % covering family
\newcommand{\Ob}{\mathrm{Ob}}              % objects
\newcommand{\Mor}{\mathrm{Mor}}            % morphisms
\newcommand{\res}{\mathrm{res}}            % restriction
\newcommand{\Cech}{\check{C}}              % Čech complex
\newcommand{\Hcoh}[1]{H^{#1}}             % cohomology group

% Descent
\newcommand{\descent}{\mathrm{desc}}
\newcommand{\glue}{\mathrm{glue}}
\newcommand{\overlap}[2]{#1 \cap #2}

% Semantic moves
\newcommand{\VERIFY}{\textsc{Verify}}
\newcommand{\CONSTRUCT}{\textsc{Construct}}
\newcommand{\NORMALIZE}{\textsc{Normalize}}
\newcommand{\GLUE}{\textsc{Glue}}
\newcommand{\RESTRICT}{\textsc{Restrict}}
\newcommand{\TRANSPORT}{\textsc{Transport}}
\newcommand{\REPAIR}{\textsc{Repair}}
\newcommand{\PROMOTE}{\textsc{Promote}}
\newcommand{\CHALLENGE}{\textsc{Challenge}}
\newcommand{\ESCALATE}{\textsc{Escalate}}

% Fragments
\newcommand{\QFLIA}{\texttt{QF\_LIA}}
\newcommand{\QFLRA}{\texttt{QF\_LRA}}
\newcommand{\QFBV}{\texttt{QF\_BV}}
\newcommand{\QFUF}{\texttt{QF\_UF}}

% Misc
\newcommand{\Python}{\textsc{Python}}
\newcommand{\JuGeo}{\textsc{JuGeo}}
\newcommand{\LEAN}{\textsc{Lean}\,4}
\newcommand{\Fstar}{F$^{\star}$}
\newcommand{\Coq}{\textsc{Coq}}
\newcommand{\Dafny}{\textsc{Dafny}}

% Common shortcuts
\newcommand{\ie}{i.e.\@\xspace}
\newcommand{\eg}{e.g.\@\xspace}
\newcommand{\cf}{cf.\@\xspace}
\newcommand{\wrt}{w.r.t.\@\xspace}

% Evaluation shortcuts
\newcommand{\numcases}{300}
\newcommand{\numspec}{100}
\newcommand{\numeq}{100}
\newcommand{\numbug}{100}
\newcommand{\medtime}{6.5\,ms}
\newcommand{\totaltime}{1.949\,s}


% -- Additional packages -------------------------------------------------
\usepackage{array}
\usepackage{tikz}
\usetikzlibrary{arrows.meta,positioning,calc,fit,backgrounds,decorations.markings,decorations.pathmorphing,shapes.geometric}
\usepackage{algorithm}
\usepackage{algorithmic}

% -- Macros --------------------------------------------------------------
\newcommand{\JB}{\mathcal{J}}
\newcommand{\Bundle}{\mathcal{E}}
\newcommand{\Base}{\mathcal{B}}
\newcommand{\Fiber}{\mathcal{F}}
\newcommand{\Trust}{\mathfrak{T}}
\newcommand{\Agents}{\mathcal{A}}
\newcommand{\Tasks}{\mathcal{Q}}
\newcommand{\Sections}{\Gamma}
\newcommand{\HCech}{\check{H}}
\newcommand{\Conn}{\nabla}
\newcommand{\Curv}{\Omega}
\newcommand{\Hol}{\mathrm{Hol}}
\newcommand{\Chern}{c_1}
\newcommand{\Claim}{\sigma}
\newcommand{\Evidence}{e}
\newcommand{\Channel}{\chi}
\newcommand{\TrustVar}{\tau}
\newcommand{\Transport}{\Delta_{\tau}}
\newcommand{\JuGeoAgents}{\texttt{jugeo-agents}}
\newcommand{\ClaudeCode}{\textsc{Claude Code}}
\newcommand{\CopilotCLI}{\textsc{Copilot CLI}}
\newcommand{\CodexCLI}{\textsc{Codex}}
\newcommand{\GlobalAssembler}{\textsc{GlobalSectionAssembler}}
\newcommand{\BundleDiag}{\textsc{JudgmentBundle}}
\newcommand{\WebBundle}{\textsc{WebTrustBundle}}
\newcommand{\TrustFlat}{\mathsf{Flat}}
\newcommand{\TrustCurved}{\mathsf{Curved}}
\newcommand{\Echo}{\mathsf{Echo}}
\newcommand{\Phantom}{\mathsf{Phantom}}
\newcommand{\RAG}{\mathsf{RAG}}
\newcommand{\ToolExec}{\mathsf{TOOL\_EXECUTED}}
\newcommand{\ToolVerified}{\mathsf{TOOL\_VERIFIED}}
\newcommand{\Citation}{\mathsf{CITATION\_BACKED}}
\newcommand{\Ungrounded}{\mathsf{UNGROUNDED\_CLAIM}}
\newcommand{\Formal}{\mathsf{FORMALLY\_PROVEN}}
\newcommand{\LLM}{\mathsf{MODEL\_GENERATED}}

\title{\textbf{Judgment Fiber Bundles: Trust as a Geometric Connection in Multi-Agent AI Verification}}
\author{JuGeo Research Group}
\date{\today}

\begin{document}
\maketitle

\begin{abstract}
Multi-agent AI systems now routinely decompose tasks across several tool-using
language models, but they still lack formal foundations for deciding when a
collection of agent outputs should be trusted as a whole.  Existing techniques
such as majority vote, debate, self-reflection, and pairwise adjudication are
fundamentally local: they compare outputs one edge at a time and therefore miss
structural team pathologies.  We introduce the \emph{judgment fiber bundle}, a
construction at the intersection of differential geometry, sheaf theory, and
formal verification.  A judgment $(\Claim, \Evidence, \TrustVar, \Channel)$ is a
claim together with its evidence, trust level, and evidence channel.  These
judgments form the fibers of a bundle $\pi : \Bundle \to \Base$ over an
agent-task site.  Trust is not treated as a flat label attached to an output;
instead, trust transport between agents defines a connection $\Conn$ on the
bundle.  Its curvature detects structural team unreliability invisible to
pairwise checks, while its holonomy detects loop-level defects such as trust
monopolies, echo chambers, and cascades.  The first Chern class $\Chern$
provides a computable global invariant: $\Chern = 0$ if and only if a globally
consistent trust assignment exists.  We show that the judgment bundle strictly
generalizes \v{C}ech cohomological consistency checking: $H^1$ contradictions
embed as localized curvature, $H^2$ cascades as holonomy, and phantom consensus
as flat-but-evidentially-empty regions.  We describe an implementation in
\JuGeoAgents{}, including adapters for \ClaudeCode{}, \CopilotCLI{}, and
\CodexCLI{}, and present coding and web-stack examples in which bundle geometry
surfaces decisive diagnostics unavailable to naive consensus.
\end{abstract}

\newpage

\section{Introduction}\label{sec:introduction}

Multi-agent AI is moving rapidly from research curiosity to engineering
infrastructure.  Contemporary systems assign different subtasks to different
language models, permit specialized tools or retrieval pathways for each agent,
and then attempt to aggregate the resulting outputs into a single decision or
artifact.  The pattern appears across coding assistants, research agents,
security analyzers, customer-support orchestrators, and autonomous browsing
systems.  Yet despite this practical importance, the core mathematical question
is unsettled: \emph{what does it mean for a team of AI agents to be globally
consistent}?  Put differently, under what conditions do local outputs produced
under different evidence conditions, prompt contexts, trust ceilings, and tool
permissions combine into a single reliable global section of work?

The dominant answers in the literature are mostly procedural.  Majority vote is
the most common baseline; more elaborate variants use self-consistency,
chain-of-thought comparison, debate, or confidence-weighted ranking.  These
methods improve robustness in some regimes, but they remain essentially
0-dimensional.  They compare endpoints without modeling the geometry of trust
between them.  If Agent~A trusts Agent~B because B used tools, and Agent~B
trusts Agent~C because C sounds coherent, then A may end up indirectly granting
C far more authority than C deserves.  This is not merely an error in scoring;
it is a failure of \emph{transport}.  The meaning of trust changes along the
path by which a judgment moves through the system.

Judgment Geometry suggests a different answer.  In the broader JuGeo program,
judgments do not live in a flat list but on semantic sites equipped with covers,
restriction maps, and descent conditions.  Obstructions to gluing are detected
cohomologically.  That perspective already improves on naive voting because it
asks whether local sections agree on overlaps.  But even that is not enough for
multi-agent systems.  Sheaf cohomology detects whether claims glue; it does not,
by itself, measure whether the \emph{trust field} used to compare those claims is
internally coherent.  We therefore introduce a differential-geometric layer on
top of the sheaf layer.

\paragraph{Trust is not a label.}
In most AI systems, trust is modeled as a scalar or categorical label attached to
an output.  We argue that this is the wrong primitive.  Trust is better modeled
as a \emph{connection}: it specifies how a judgment's trust level changes when
that judgment is moved from one agent-task context to another.  Once trust is
modeled as a connection, familiar geometric notions become available.  Curvature
measures path dependence.  Holonomy measures the total trust defect around a
loop.  Characteristic classes summarize global structure.  Echo chambers,
phantom consensus, cascades, and monopolies become geometric phenomena rather
than heuristics.

The advantage is practical as well as conceptual.  Bundle diagnostics tell us
not only whether a team disagrees, but \emph{what kind of team it is}.  A team
with $\Chern > 0$ exhibits trust inflation: mutually reinforcing weak evidence
creates unjustified certainty.  A team with $\Chern < 0$ exhibits trust
deflation: strong evidence is repeatedly attenuated across boundaries, producing
unnecessary conservatism or hidden single-agent dominance.  A team with trivial
pairwise disagreements but non-trivial holonomy is structurally unsound even if
all local tests appear fine.

This paper makes five contributions.
\begin{enumerate}
  \item We define the \emph{judgment fiber bundle}, whose fibers consist of
    judgments $(\Claim, \Evidence, \TrustVar, \Channel)$ above an agent-task site
    (\cref{sec:bundle-construction}).
  \item We define trust transport as a connection and prove discrete analogues
    of gauge invariance and the Ambrose--Singer relation between holonomy and
    curvature (\cref{sec:bundle-construction,sec:characteristic-classes}).
  \item We show that JuGeo's existing \v{C}ech cohomological machinery embeds
    into the bundle framework, with $H^1$ contradictions corresponding to local
    curvature and $H^2$ cascades corresponding to holonomy
    (\cref{sec:cech-embedding}).
  \item We describe a concrete implementation in \JuGeoAgents{}, including
    adapters for \ClaudeCode{}, \CopilotCLI{}, and \CodexCLI{}, together with
    web-stack trust bundles and cross-language bundle diagnostics
    (\cref{sec:implementation}).
  \item We provide worked examples showing how bundle geometry rapidly exposes
    useful structure in coding and web verification tasks
    (\cref{sec:examples,sec:webapps}).
\end{enumerate}

The paper is self-contained, but it should be read as part of the JuGeo series:
semantic sites provide the base category, trust algebra provides the ordered
epistemic lattice, and descent provides the sheaf-theoretic gluing criterion.
The judgment bundle adds a new statement: even when gluing data are available,
there may still be a geometric obstruction in the trust field itself.

\section{Preliminaries}\label{sec:preliminaries}

\subsection{Fiber bundles and connections}

A fiber bundle consists of a total space $\Bundle$, a base space $\Base$, a
projection map $\pi : \Bundle \to \Base$, and fibers $\Fiber_p = \pi^{-1}(p)$
that are locally but not necessarily globally trivial.  In our setting, the
base space is discrete and combinatorial rather than smooth, but the bundle
intuition remains useful: each point in the base determines a context, and the
fiber above that point contains the judgments available in that context.

A connection $\Conn$ specifies how to compare neighboring fibers.  In the smooth
setting it can be viewed as a horizontal sub-bundle or a covariant derivative;
in the discrete setting used here it is more naturally a transport rule on edges
of the base nerve.  Given neighboring points $p,q \in \Base$, the connection
provides a parallel transport map $\Conn_{p \to q} : \Fiber_p \to \Fiber_q$ that
adjusts trust according to the evidential relation between contexts.  Curvature
measures the failure of transport to be path-independent.  Holonomy measures the
net effect of transport around a closed loop.

\begin{figure}[H]
\centering
\begin{tikzpicture}[scale=1.15]
  \draw[thick, jugeoBlue] (0,0) -- (6,0);
  \node[below] at (3,0) {$\Base$ (task space)};

  \foreach \x/\label in {1/A,3/B,5/C} {
    \draw[thick, jugeoPurple] (\x,0) -- (\x,3);
    \node[below, jugeoBlue] at (\x,0) {$\label$};
    \node[above, jugeoPurple] at (\x,3) {$\Fiber_{\label}$};
    \fill[jugeoGreen] (\x,0.7) circle (2pt);
    \fill[jugeoGreen] (\x,1.55) circle (2pt);
    \fill[jugeoRed] (\x,2.25) circle (2pt);
  }

  \draw[-{Stealth}, thick, orange, dashed] (1.1,1.5) .. controls (2,2.1) .. (2.9,1.8);
  \node[above, orange, font=\small] at (2,2.25) {$\Conn_{A\to B}$};
  \draw[-{Stealth}, thick, orange, dashed] (3.1,1.8) .. controls (4,2.25) .. (4.9,2.15);
  \node[above, orange, font=\small] at (4,2.45) {$\Conn_{B\to C}$};

  \node[right] at (6.2,1.6) {$\Bundle$ (judgment space)};
\end{tikzpicture}
\caption{A discrete fiber bundle view of a multi-agent task: points in the base
are agent-task contexts; fibers contain judgments in those contexts; dashed
arrows represent trust transport.}
\label{fig:fiber-schematic}
\end{figure}

The first Chern class is, in the smooth complex-line-bundle setting, the de Rham
cohomology class represented by curvature.  In our discrete setting we use an
average-curvature surrogate that plays the same practical role: it summarizes
whether a globally flat trust assignment exists.  This is not merely notation
borrowing.  The essential phenomenon is the same: local transport data can carry
a global obstruction.

\subsection{\v{C}ech cohomology and descent}

Let $\mathcal{U} = \{U_i\}$ be a cover of a semantic site and let $\mathscr{S}$
be a sheaf of judgments or proposals on that site.  A 0-cochain is a family of
local sections $\{s_i \in \mathscr{S}(U_i)\}$.  Its coboundary measures pairwise
mismatch on overlaps $U_i \cap U_j$.  A vanishing coboundary means the local
sections satisfy the cocycle condition.  When those compatible local sections
glue to a global section, descent succeeds.  When they do not, the obstruction
appears in $\HCech^1$ and its cohomology class records the failure.

JuGeo uses this structure operationally.  Agent outputs are turned into local
sections; overlaps are computed from task decompositions, shared claims, or
cross-language contracts; contradiction detection yields $H^1$-style
obstructions; cascades and phantom consensus behave like higher obstructions.
The present paper does not replace that sheaf machinery.  It asks a second
question: how coherent is the trust structure used to compare those local
sections?

\subsection{Trust algebra in Judgment Geometry}

JuGeo does not treat trust as a scalar confidence score.  It is an ordered
algebra with explicit ceilings, promotion rules, demotion rules, and
channel-sensitive constraints.  In \JuGeoAgents{}, the canonical lattice is
ordered as follows:
\[
  \mathsf{SELF\_CONTRADICTED} < \Ungrounded < \mathsf{WEAK\_MODEL\_GENERATED}
  < \mathsf{STRONG\_MODEL\_GENERATED} < \Citation < \RAG < \ToolExec
  < \ToolVerified < \mathsf{CROSS\_AGENT\_CONFIRMED} < \Formal.
\]

Three axioms matter for the present paper.
\begin{enumerate}
  \item \textbf{No silent promotion.}  Trust may increase only via explicit,
    named evidence routes.
  \item \textbf{Conservative join.}  Aggregating evidence cannot silently exceed
    the strongest admissible supporting channel.
  \item \textbf{Evidence-gated promotion.}  Each channel has a ceiling: a pure
    model response cannot become tool-verified without tool evidence.
\end{enumerate}

These axioms already make trust algebraic.  The judgment bundle makes it
geometric by studying how trust changes under transport.

\section{The judgment fiber bundle}\label{sec:bundle-construction}

\subsection{The base space: the agent-task site}

\begin{definition}[Judgment]
A \emph{judgment} is a tuple
\[
  J = (\Claim, \Evidence, \TrustVar, \Channel)
\]
where $\Claim$ is a proposition or claim, $\Evidence$ is a finite evidence set,
$\TrustVar \in \Trust$ is a trust level, and $\Channel$ is the channel through
which the evidence was obtained.
\end{definition}

\begin{definition}[Agent-task site]
Let $\Agents = \{a_1,\dots,a_m\}$ be a set of agents and
$\Tasks = \{q_1,\dots,q_n\}$ a task decomposition.  The \emph{agent-task site}
$\Base$ is the nerve of the covering induced by assigning subtasks to agents.
Its points are pairs $p=(a_i,q_j)$, and its edges connect points that overlap on
shared specifications, files, interfaces, or claims.  Higher simplices encode
multi-way overlap.
\end{definition}

\begin{figure}[H]
\centering
\begin{tikzpicture}[scale=1.0]
  \node[circle, draw, fill=jugeoBlue!20, minimum size=1.2cm] (A) at (0,0) {Agent A};
  \node[circle, draw, fill=jugeoBlue!20, minimum size=1.2cm] (B) at (4,0) {Agent B};
  \node[circle, draw, fill=jugeoBlue!20, minimum size=1.2cm] (C) at (2,3) {Agent C};

  \draw[thick] (A) -- (B) node[midway, below] {overlap};
  \draw[thick] (B) -- (C) node[midway, right] {overlap};
  \draw[thick] (C) -- (A) node[midway, left] {overlap};

  \node at (2,1.2) {$\sigma_{ABC}$};
  \node[red, font=\bfseries] at (2,0.35) {$F = \Delta\tau_{AB}+\Delta\tau_{BC}+\Delta\tau_{CA}$};
\end{tikzpicture}
\caption{The nerve of an agent-task cover.  Curvature is computed on the
2-simplex induced by three mutually overlapping agents.}
\label{fig:agent-nerve}
\end{figure}

\subsection{Fibers and total space}

\begin{definition}[Judgment fiber]
For each point $p \in \Base$, the \emph{judgment fiber} above $p$ is
\[
  \Fiber_p = \{ (\Claim, \Evidence, \TrustVar, \Channel) :
  \text{the context } p \text{ produces a judgment about } \Claim \}.
\]
\end{definition}

\begin{definition}[Judgment bundle]
The \emph{judgment bundle} is the projection
\[
  \pi : \Bundle \to \Base,
\]
where the total space is the disjoint union
\[
  \Bundle = \coprod_{p \in \Base} \Fiber_p.
\]
The structure group is the trust lattice acting by admissible transport on the
trust coordinate of each judgment.
\end{definition}

Intuitively, the bundle separates two kinds of structure that are often
confounded in agent systems.  The base remembers \emph{where} a judgment came
from: which agent, subtask, interface, or lifecycle stage.  The fiber remembers
\emph{what} was judged and under what evidential conditions.  This separation is
what makes trust transport meaningful: trust changes not because the claim has
changed, but because the context in which it is compared has changed.

\begin{equation}
\begin{tikzcd}
\Fiber_A \arrow[r, "\Conn_{A\to B}"] \arrow[d, "\pi"'] & \Fiber_B \arrow[d, "\pi"] \\
A \arrow[r] & B
\end{tikzcd}
\end{equation}

\subsection{Trust as a connection}

\begin{definition}[Trust connection]
A \emph{trust connection} on the judgment bundle assigns to each oriented edge
$p \to q$ in the nerve of $\Base$ a transport map
\[
  \Conn_{p\to q} : \Fiber_p \to \Fiber_q
\]
that acts by transporting judgments with the same or overlapping semantic
content while adjusting trust according to the evidential relationship between
$p$ and $q$.
\end{definition}

Operationally, \JuGeoAgents{} constructs this connection from observed trust
pairs.  If Agent~A makes a tool-verified claim and Agent~B makes an otherwise
matching citation-backed claim, then transport from A to B carries a negative
trust delta.  If a lower-trust agent's claim becomes stronger only after a tool
check at the target, then transport may carry a positive delta.  When no
relevant comparison is available, the transport is taken to be flat.

\begin{definition}[Curvature]
For a 2-simplex $(A,B,C)$ in the base nerve, the curvature is
\[
  \Curv(A,B,C) = \Transport(A\to B) + \Transport(B\to C) + \Transport(C\to A).
\]
Non-zero curvature means trust transport around the triangle is path-dependent.
\end{definition}

\begin{theorem}[Gauge invariance]
Let all trust levels in a connected region be shifted by a constant $c$ in the
strength index representation.  Then the curvature of every 2-simplex is
unchanged.
\end{theorem}

\begin{proof}
If the transported trust differential on an edge is
$\Transport(A\to B)=\TrustVar_B-\TrustVar_A$, then after the gauge shift we have
\[
  \Transport'(A\to B) = (\TrustVar_B+c) - (\TrustVar_A+c) = \TrustVar_B-\TrustVar_A.
\]
Therefore each edge differential is unchanged, and hence their sum around any
triangle is unchanged as well.
\end{proof}

\begin{theorem}[Discrete Ambrose--Singer]
Let $\gamma$ be a closed loop in the 1-skeleton of the base nerve, and let
$\mathcal{T}(\gamma)$ be any triangulation of the region it bounds.  Then
\[
  \Hol(\gamma) = \sum_{\sigma \in \mathcal{T}(\gamma)} \Curv(\sigma).
\]
\end{theorem}

\begin{proof}
Write each face curvature as the oriented sum of its three edge transport
values.  Summing over all faces in the triangulation causes every interior edge
to appear exactly twice with opposite orientation, so all interior contributions
cancel.  The only surviving terms are the boundary edges, which are precisely
the transport values along the loop $\gamma$.  Their sum is the holonomy.
\end{proof}

\begin{figure}[H]
\centering
\begin{tikzpicture}[scale=1.35]
  \node[circle, draw, thick, fill=jugeoBlue!15, inner sep=3pt] (A) at (0,0) {$A$};
  \node[circle, draw, thick, fill=jugeoGreen!15, inner sep=3pt] (B) at (3,0) {$B$};
  \node[circle, draw, thick, fill=jugeoPurple!15, inner sep=3pt] (C) at (1.5,2.5) {$C$};

  \node[below left, font=\small] at (A) {$\tau=7$};
  \node[below right, font=\small] at (B) {$\tau=5$};
  \node[above, font=\small] at (C) {$\tau=3$};

  \draw[-{Stealth[length=8pt]}, thick, jugeoBlue,
    decoration={markings, mark=at position 0.5 with {\node[below, font=\small\bfseries, jugeoBlue] {$-2$};}},
    postaction=decorate] (A) -- (B);
  \draw[-{Stealth[length=8pt]}, thick, jugeoGreen,
    decoration={markings, mark=at position 0.5 with {\node[right, font=\small\bfseries, jugeoGreen] {$-2$};}},
    postaction=decorate] (B) -- (C);
  \draw[-{Stealth[length=8pt]}, thick, jugeoPurple,
    decoration={markings, mark=at position 0.5 with {\node[left, font=\small\bfseries, jugeoPurple] {$+4$};}},
    postaction=decorate] (C) -- (A);

  \node[draw, rounded corners, fill=jugeoRed!10, inner sep=4pt, font=\bfseries] at (1.5,0.8)
    {$\Curv = -2 + (-2) + 4 = 0$};
\end{tikzpicture}
\caption{Holonomy around a three-agent loop.  Even in a small team, trust
transport may or may not close consistently.}
\label{fig:holonomy-loop}
\end{figure}

\section{Characteristic classes and structural pathologies}\label{sec:characteristic-classes}

\begin{definition}[First Chern class]
Let $\mathcal{F}_2$ denote the set of 2-faces in the base nerve.  The discrete
first Chern class of the judgment bundle is the average curvature
\[
  \Chern = \frac{1}{|\mathcal{F}_2|} \sum_{\sigma \in \mathcal{F}_2} \Curv(\sigma).
\]
\end{definition}

\begin{theorem}[Flatness criterion]
Suppose every connected component of the base nerve is triangulated and every
edge transport is determined by admissible trust differentials.  Then
$\Chern = 0$ and trivial loop holonomy if and only if the judgment bundle admits
a globally consistent flat trust assignment on that component.
\end{theorem}

\begin{proof}
If the bundle admits a globally flat assignment, then every loop has trivial
holonomy and every triangular face has zero curvature; therefore the average
curvature is zero.  Conversely, if the average curvature is zero and every loop
holonomy is trivial, then edge transport is path-independent on the connected
component.  Fix a basepoint and assign its trust arbitrarily; path independence
extends this assignment uniquely to every other point.  The resulting global
assignment realizes the observed transports without contradiction, so the
connection is flat.
\end{proof}

\begin{corollary}
If $\Chern \neq 0$, then there exists no globally consistent trust assignment for
the team.  The unreliability is structural rather than local.
\end{corollary}

The sign of $\Chern$ carries actionable meaning.
\begin{center}
\begin{tabular}{@{}lll@{}}
\toprule
Value of $\Chern$ & Interpretation & Team pathology \\
\midrule
$\Chern = 0$ & Flat bundle & Healthy team geometry \\
$\Chern > 0$ & Trust inflation & Echo chamber or circular citation \\
$\Chern < 0$ & Trust deflation & Over-attenuation or hidden monopoly \\
Large $|\Chern|$ & High curvature & Severe structural instability \\
\bottomrule
\end{tabular}
\end{center}

These interpretations are visible in practice.  Positive curvature indicates
that trust increases as judgments circulate, usually because weak agents cite or
echo each other until the team behaves as though independent confirmation has
occurred.  Negative curvature indicates repeated attenuation, often because one
agent really has the only strong evidence while the others merely relay weaker
versions of it.  In both cases, the bundle gives a team-level diagnosis rather
than merely ranking individuals.

\section{Embedding \v{C}ech cohomology into bundle geometry}\label{sec:cech-embedding}

The bundle framework is designed to sit on top of JuGeo's existing sheaf layer,
not to replace it.  The crucial relationship is that the failure of sections to
glue in the sheaf appears as non-flat geometry in the bundle.

\begin{theorem}[Embedding theorem]
There is a natural transformation from the obstruction calculus of the judgment
sheaf to the curvature--holonomy calculus of the judgment bundle such that:
\begin{enumerate}
  \item $H^0$ global sections correspond to global flat sections of the bundle;
  \item $H^1$ contradictions map to localized non-zero curvature on faces
    containing the contradictory overlap;
  \item $H^2$ cascades map to non-trivial holonomy around loops induced by the
    cascade structure;
  \item phantom consensus maps to flat local transport with evidentially empty
    or ungrounded fibers.
\end{enumerate}
\end{theorem}

\begin{proof}
A compatible family of local sections defines both a sheaf-theoretic cocycle and
a family of judgments in overlapping fibers.  If the family glues globally, the
transport values induced by matching judgments compose consistently, so the
associated bundle section is flat.  If a contradiction appears on an overlap,
then the induced transport around any face spanning that overlap cannot sum to
zero, yielding curvature.  Higher-order incompatibilities that evade pairwise
checks require traversing loops in the overlap structure and therefore manifest
as holonomy.  Finally, phantom consensus consists of globally compatible but
ungrounded claims; these induce flat transport because no explicit contradiction
is present, yet the fibers remain epistemically empty.
\end{proof}

\begin{equation}
\begin{tikzcd}
\HCech^1(\mathcal{U},\JB) \arrow[r, hook, "\iota"] \arrow[d, "\text{detects}"'] & \Omega^2(\Bundle/\Base) \arrow[d, "\text{detects}"] \\
\text{contradictions} \arrow[r, equal] & \text{non-zero curvature}
\end{tikzcd}
\end{equation}

\begin{proposition}[Strict generalization]
There exist structurally unhealthy multi-agent teams with trivial local
contradiction sets but non-zero characteristic data.  In particular, echo
chambers and trust monopolies may be invisible to the raw \v{C}ech obstruction
set while remaining visible through $\Chern$ or non-trivial holonomy.
\end{proposition}

\begin{proof}
Consider three agents that all repeat the same claim, so there is no pairwise
contradiction and thus no $H^1$ witness.  Suppose, however, that the claim is
ungrounded and each agent's trust is promoted only because the others repeat it.
Then transport around the triangle inflates trust, producing positive curvature
without an explicit contradiction.  Likewise, if one agent has all the strong
evidence and the others merely echo or attenuate it, pairwise consistency can
still hold while loop transport records a non-trivial monopoly pattern.
\end{proof}

The practical dictionary is therefore:
\[
\begin{aligned}
H^0 &\leftrightarrow \text{verified global sections},\\
H^1 &\leftrightarrow \text{curvature hotspots},\\
H^2 &\leftrightarrow \text{holonomy defects},\\
\text{phantom classes} &\leftrightarrow \text{flat but ungrounded fibers}.
\end{aligned}
\]

\section{Implementation in \JuGeoAgents{}}\label{sec:implementation}

\subsection{Architecture}

The implementation is distributed across three layers.  The shared schema layer
contains trust levels, factual claims, contradictions, provenance structures,
and agent-output dataclasses.  The sheaf/cohomology layer contains claim
extraction, contradiction detection, descent checking, and the global-section
assembler.  The bundle layer adds trust transport, curvature, holonomy,
stratification, and characteristic classes.  Adapters lift external systems into
that common runtime.

\begin{figure}[H]
\centering
\begin{tikzpicture}[
  box/.style={draw, rounded corners, minimum width=3cm, minimum height=0.8cm, font=\small\bfseries},
  agent/.style={box, fill=jugeoBlue!20},
  core/.style={box, fill=jugeoGreen!20},
  output/.style={box, fill=jugeoPurple!20},
]
  \node[agent] (claude) at (0,4) {Claude Code};
  \node[agent] (copilot) at (4,4) {Copilot CLI};
  \node[agent] (codex) at (8,4) {Codex};

  \node[box, fill=orange!20, minimum width=9cm] (adapters) at (4,3) {Coding Agent Adapters};

  \node[core] (bundle) at (0,1.5) {JudgmentBundle};
  \node[core] (fusion) at (4,1.5) {GlobalSectionAssembler};
  \node[core] (descent) at (8,1.5) {DescentEngine};

  \node[output, minimum width=9cm] (report) at (4,0) {PipelineReport + Bundle Diagnostics};

  \draw[-{Stealth}] (claude) -- (adapters);
  \draw[-{Stealth}] (copilot) -- (adapters);
  \draw[-{Stealth}] (codex) -- (adapters);
  \draw[-{Stealth}] (adapters) -- (bundle);
  \draw[-{Stealth}] (adapters) -- (fusion);
  \draw[-{Stealth}] (adapters) -- (descent);
  \draw[-{Stealth}] (bundle) -- (report);
  \draw[-{Stealth}] (fusion) -- (report);
  \draw[-{Stealth}] (descent) -- (report);
\end{tikzpicture}
\caption{Implementation architecture in \JuGeoAgents{}.  External coding agents
are normalized into a common runtime, then analyzed simultaneously by descent,
cohomology, and bundle geometry.}
\label{fig:architecture}
\end{figure}

The recent retrofit performed across the JuGeo ecosystem realizes this design in
three places.
\begin{enumerate}
  \item In \texttt{jugeo-agents}, the wrapper and fusion pipelines now ingest
    every \texttt{AgentOutput} into \BundleDiag{}, and bundle diagnostics are
    attached to pipeline reports.
  \item In the JuGeo core library, \texttt{jugeo.geometry.bundle} introduces a
    site-native \texttt{VerificationBundle} so the same ideas apply to semantic
    coordinates, proofs, tests, and theorem obligations.
  \item In the web stack, \texttt{jugeo.webapp.trust.bundle},
    \texttt{FiberedBundleDiagnostics}, and \texttt{DeviceBundleDiagnostics}
    bring connection geometry to request lifecycles, language boundaries, and
    cross-device visual invariants.
\end{enumerate}

\subsection{A coding-agents example: Claude Code, Copilot CLI, Codex}

The most intuitive use-case is competitive coding.  Three agents receive the
same specification: implement a thread-safe token-bucket rate limiter in
Python.  They produce code under different evidence conditions.
\ClaudeCode{} uses tools to run tests.  \CopilotCLI{} cites a plausible
reference implementation but does not execute anything.  \CodexCLI{} produces a
reasonable decorator from pure generation alone.

\begin{lstlisting}[style=jugeo-python,caption={Competitive coding with three agent CLIs.},label={lst:coding-agents}]
from jugeo_agents import (
    JudgmentBundle,
    GlobalSectionAssembler,
    AgentOutput,
    CodingAgentOrchestrator,
)

orch = CodingAgentOrchestrator(
    task="Implement a thread-safe token bucket rate limiter in Python"
)

claude_output = orch.submit_agent_output(AgentOutput(
    agent_id="claude-code",
    model="claude-sonnet-4",
    output_text="""
import threading
import time

class TokenBucket:
    def __init__(self, rate: float, capacity: int):
        self.rate = rate
        self.capacity = capacity
        self._tokens = capacity
        self._last_refill = time.monotonic()
        self._lock = threading.Lock()

    def _refill(self) -> None:
        now = time.monotonic()
        elapsed = now - self._last_refill
        self._tokens = min(self.capacity,
                           self._tokens + elapsed * self.rate)
        self._last_refill = now

    def consume(self, tokens: int = 1) -> bool:
        with self._lock:
            self._refill()
            if self._tokens >= tokens:
                self._tokens -= tokens
                return True
            return False
    """,
    tools_used=["file_write", "bash_run", "test_execute"],
    tool_results={"test_execute": "All 12 tests passed"},
))

copilot_output = orch.submit_agent_output(AgentOutput(
    agent_id="copilot-cli",
    model="gpt-4.1",
    output_text="""
import time

class RateLimiter:
    def __init__(self, max_requests, window_seconds):
        self.max_requests = max_requests
        self.window_seconds = window_seconds
        self.requests = []

    def allow(self):
        now = time.time()
        self.requests = [t for t in self.requests
                         if now - t < self.window_seconds]
        if len(self.requests) < self.max_requests:
            self.requests.append(now)
            return True
        return False
    """,
    citations=["https://example.com/token-bucket-reference"],
))

codex_output = orch.submit_agent_output(AgentOutput(
    agent_id="codex",
    model="codex-mini",
    output_text="""
import time

def rate_limit(func, max_calls=100, period=60):
    calls = []
    def wrapper(*args, **kwargs):
        now = time.time()
        calls[:] = [c for c in calls if now - c < period]
        if len(calls) >= max_calls:
            raise Exception("Rate limited")
        calls.append(now)
        return func(*args, **kwargs)
    return wrapper
    """,
))

result = orch.verify()
print(result.bundle_summary)
\end{lstlisting}

A representative bundle report is:
\begin{lstlisting}[style=jugeo-python,caption={Representative bundle diagnostic from the coding-agents example.}]
=== Judgment Fiber Bundle Diagnostic ===
  Agents: claude-code, codex, copilot-cli
  Total judgments: 23

  Connection Geometry:
    First Chern class c1 = -1.3333
    Trust deflation: tool-verified evidence systematically dominates.
    Curved faces: 1 / 1
    Bundle is flat: No

  Team Holonomy:
    Loop: claude-code -> codex -> copilot-cli
    Total shift: -8.00
    Non-trivial holonomy: hidden trust monopoly by claude-code

  Trust Stratification:
    UNGROUNDED_CLAIM: 8 judgments
    CITATION_BACKED: 7 judgments
    TOOL_VERIFIED: 8 judgments

  Recommendation: trust claude-code's implementation.
\end{lstlisting}

Two things are important here.  First, the decision is not based on a vague
sense that one model is ``better.''  It is based on evidence channels and the
resulting transport geometry.  Second, the diagnosis is more expressive than a
winner-takes-all leaderboard.  The bundle tells us that the team is not a flat
consensus but a negatively curved configuration dominated by one strong source.
That makes the right operational response obvious: trust the winning code,
quarantine the weak variants, and avoid pretending that a broad consensus has
been established.

\subsection{Knowledge fusion example}

Bundle geometry also works for research-style teams that produce natural-language
claims rather than code.  Here the bundle is paired with sheaf-based claim
extraction and the global section assembler.

\begin{lstlisting}[style=jugeo-python,caption={Knowledge fusion with cohomology plus bundle diagnostics.},label={lst:fusion-example}]
from jugeo_agents import GlobalSectionAssembler

assembler = GlobalSectionAssembler()
for output in research_team_outputs:
    assembler.ingest(output)

result = assembler.assemble()
print("verified:", len(result.verified_claims))
print("quarantined:", result.quarantined_count)
print("metadata keys:", sorted(result.metadata.keys()))
print("bundle flat:", result.metadata["bundle_diagnostics"]["bundle_is_flat"])
\end{lstlisting}

The resulting object carries both sheaf data and geometric data.  This matters
because some failure modes are primarily logical while others are primarily
geometric.  A direct contradiction about the number of qubits in a hardware
announcement is a sheaf-level issue first.  A collection of ungrounded agents
that all repeat the same market-size claim is geometric first: it forms a trust
shape before it forms a contradiction.

\subsection{A web-stack example}

The web-specific bundles introduced in the retrofit make the same ideas useful
far beyond agent-to-agent reasoning.  In a full-stack Flask-style system, the
request lifecycle
\[
  \text{browser} \to \text{route} \to \text{handler} \to \text{database}
  \to \text{response} \to \text{render}
\]
acts as a base path, while judgments about validation, sanitization, typing,
authentication, and rendering populate the fibers above each stage.

\begin{lstlisting}[style=jugeo-python,caption={Web trust bundle example.},label={lst:web-bundle-example}]
from jugeo.webapp.trust.bundle import WebTrustBundle, WebJudgment

bundle = WebTrustBundle()
bundle.add_judgment(WebJudgment(
    stage="browser",
    language="javascript",
    claim="email field validates format",
    trust="CLIENT_VALIDATED",
))
bundle.add_judgment(WebJudgment(
    stage="handler",
    language="python",
    claim="email field validates format",
    trust="SERVER_VALIDATED",
))
bundle.add_judgment(WebJudgment(
    stage="database",
    language="sql",
    claim="email column has NOT NULL constraint",
    trust="DB_CONSTRAINT_ENFORCED",
))
print(bundle.summary_text())
\end{lstlisting}

This generalization is not cosmetic.  The same geometry that diagnoses a team of
coding agents also diagnoses a cross-language web stack.  That is exactly the
kind of JuGeo unification we want: one theory of bundle transport that applies
from semantic proofs to Flask handlers to CLI coding agents.

\section{Worked examples}\label{sec:examples}

\subsection{Echo chamber detection}

Consider three agents A, B, and C.  None has tool access.  Each cites the other
two as supporting evidence.  The textual claims match, so naive majority vote
accepts the consensus.  Pairwise contradiction detection also remains silent.
Yet the consensus is structurally illegitimate because all evidence is circular.
In the bundle this appears as positive curvature.

\begin{figure}[H]
\centering
\begin{tikzpicture}[scale=1.3]
  \node[circle, draw, thick, fill=jugeoRed!20, minimum size=1cm] (A) at (0,0) {$A$};
  \node[circle, draw, thick, fill=jugeoRed!20, minimum size=1cm] (B) at (3,0) {$B$};
  \node[circle, draw, thick, fill=jugeoRed!20, minimum size=1cm] (C) at (1.5,2.5) {$C$};

  \draw[-{Stealth}, thick, jugeoRed, bend left=15] (A) to node[below, font=\scriptsize] {cites} (B);
  \draw[-{Stealth}, thick, jugeoRed, bend left=15] (B) to node[right, font=\scriptsize] {cites} (C);
  \draw[-{Stealth}, thick, jugeoRed, bend left=15] (C) to node[left, font=\scriptsize] {cites} (A);

  \node[draw, rounded corners, fill=jugeoRed!10, font=\small\bfseries] at (1.5,0.8) {$\Chern > 0$: echo chamber};
  \node[font=\scriptsize, jugeoGray] at (1.5,-0.55) {No agent has tool or solver evidence};
\end{tikzpicture}
\caption{Positive curvature as echo chamber: apparent consensus without any
independent grounding.}
\label{fig:echo-chamber}
\end{figure}

The appropriate intervention is not ``pick the best wording'' but ``demand a new
evidence channel.''  Curvature tells the orchestrator to route the disputed claim
to tools, retrieval, or formal proof rather than to more debate.

\subsection{Trust monopoly detection}

Now suppose one agent really does have the only strong evidence.  The other two
agree with it, but only after attenuation through weaker channels.  Pairwise
agreement looks reassuring.  The bundle, however, records a negative holonomy:
transport around the team loses trust in a way that reveals the team's effective
rank structure.  This is not an echo chamber but a disguised single-expert
configuration.  Again, the bundle helps the orchestrator respond correctly:
trust the strong agent, but do not misreport the result as broad consensus.

\subsection{Cross-language web verification}\label{sec:webapps}

The same geometry becomes particularly compelling in web applications because the
most expensive bugs live in overlaps between layers rather than within any single
language.  Frontend code may validate a field that the backend fails to check;
a Flask handler may assume non-null database values that the schema does not
enforce; HTML may assume an authorization state that middleware never proves.
Judgment Geometry treats these as descent failures.  Bundle geometry refines the
picture by measuring the trust transport across the browser--server--database
loop.

\begin{figure}[H]
\centering
\begin{tikzpicture}[scale=0.95]
  \foreach \i/\name/\trust/\lang in {
    0/Browser/3/JS,
    2/Route/6/Python,
    4/Handler/7/Python,
    6/Database/8/SQL,
    8/Response/5/Python,
    10/Render/3/JS
  } {
    \draw[thick, jugeoPurple] (\i,0) -- (\i,2.5);
    \node[below, font=\small\bfseries] at (\i,-0.2) {\name};
    \fill[jugeoGreen] (\i,\trust*0.3) circle (3pt);
    \node[right, font=\tiny] at (\i+0.1,\trust*0.3) {$\tau=\trust$};
    \node[above, font=\tiny, jugeoBlue] at (\i,2.5) {\lang};
  }

  \draw[thick, jugeoBlue] (-0.5,0) -- (10.5,0);
  \foreach \x in {0,2,4,6,8} {
    \pgfmathsetmacro{\nx}{\x+2}
    \draw[-{Stealth}, dashed, orange] (\x+0.2,1.5) -- (\nx-0.2,1.5);
  }
  \draw[-{Stealth}, thick, jugeoRed, bend left=30] (10.3,1) to node[above, font=\scriptsize] {$\Hol(\gamma)=-5$} (-0.3,1);
\end{tikzpicture}
\caption{A web request lifecycle as a trust bundle.  Cross-language transitions
cause trust loss that accumulates around the lifecycle loop.}
\label{fig:web-bundle}
\end{figure}

When the web bundle reports non-trivial holonomy, the message is immediate and
useful: trust is leaking across boundaries.  The system may be locally tested at
each stage, yet still globally unreliable because the evidence does not transport
cleanly through the full loop.

\section{Trust stratification and the action principle}\label{sec:stratification}

Bundle geometry naturally induces a stratification by trust level.  Judgments at
$\ToolVerified$ form a different epistemic stratum from judgments at
$\Citation$ or $\Ungrounded$.  This matters because contradictions within a
stratum are more severe than contradictions across strata.  If two tool-verified
agents disagree, there is no easy trust-based tiebreaker.  If a tool-verified
agent disagrees with an ungrounded one, the geometry is already telling us which
stratum should dominate.

\begin{figure}[H]
\centering
\begin{tikzpicture}[x=1.2cm,y=0.8cm]
  \draw[fill=jugeoRed!8] (0,0) rectangle (8,1);
  \draw[fill=orange!12] (0,1.2) rectangle (8,2.2);
  \draw[fill=jugeoBlue!10] (0,2.4) rectangle (8,3.4);
  \draw[fill=jugeoGreen!10] (0,3.6) rectangle (8,4.6);

  \node[left] at (0,0.5) {$\Ungrounded$};
  \node[left] at (0,1.7) {$\Citation$};
  \node[left] at (0,2.9) {$\ToolExec$};
  \node[left] at (0,4.1) {$\ToolVerified$};

  \foreach \x/\y in {1/0.5,2.2/0.5,4.1/0.5,6/0.5,1.5/1.7,3.4/1.7,5.6/1.7,6.9/1.7,1.2/2.9,3/2.9,5.4/2.9,7/2.9,2/4.1,4.3/4.1,6.4/4.1} {
    \fill[jugeoPurple] (\x,\y) circle (2.2pt);
  }

  \draw[thick, jugeoGray, dashed] (0.5,2.2) -- (7.5,2.4);
  \draw[thick, jugeoGray, dashed] (0.5,3.4) -- (7.5,3.6);
\end{tikzpicture}
\caption{Trust stratification of the judgment bundle.  Higher strata dominate
lower ones in fusion, while intra-stratum contradictions are maximally serious.}
\label{fig:stratification}
\end{figure}

This stratification leads naturally to an action principle for verification.
A \emph{semantic move} is any operation that changes the bundle: challenge an
agent, run a tool, request a citation, launch a proof search, or synthesize a
treaty.  Let the action of a move be
\[
  S(m) = \mathrm{cost}(m) - \Delta \Trust(m),
\]
where $\Delta \Trust(m)$ is the trust improvement produced by the move.
Verification then becomes a geodesic problem: choose the path of moves that
minimizes total action while reaching the desired trust stratum.  This gives
Judgment Geometry a precise answer to a practical orchestration question: not
just \emph{what is wrong}, but \emph{what is the cheapest principled next step}?

\section{Why this helps ideate better systems}\label{sec:ideation}

One of the surprising outcomes of the bundle perspective is that it helps with
\emph{design}, not only verification.  If a proposed application architecture
induces pathological curvature before implementation, then the design is already
suspicious.  This observation is especially powerful in full-stack and
multi-agent settings.

For coding-agent systems, bundle-aware ideation favors architectures in which
strong channels are intentionally distributed rather than concentrated.  A team
of three nearly identical weak agents is geometrically boring and often
useless.  A team in which one agent owns all tool access and the others only
summarize it may work, but it forms a monopoly geometry and therefore offers
little redundancy.  The most promising designs are those in which different
agents possess different strong channels whose overlaps can actually corroborate
one another.

For web applications, bundle-aware ideation suggests looking for app ideas whose
value comes from maintaining invariants across boundaries.  A mundane CRUD app
has little geometric interest.  By contrast, a Flask system that continuously
verifies visually stable public dashboards across device classes, or that tracks
policy conformance across browser, server, and database layers, has rich
boundary structure and therefore benefits directly from JuGeo.  This is one path
toward novel useful Flask apps: target domains where correctness lives in the
intersections and where the product itself can expose its own descent and bundle
diagnostics to users.

\section{Related work}\label{sec:related}

The closest adjacent work comes from several directions, none of which provides
a comparable geometric account.  Debate-style AI safety protocols
\cite{irving2018debate} use adversarial comparison to improve truthfulness, but
do not model trust transport.  Reflection and self-consistency methods improve
local reasoning, but again remain pointwise.  Multi-agent orchestration systems
such as AutoGen \cite{wu2023autogen}, MetaGPT \cite{hong2023metagpt}, and
CrewAI \cite{crewai2024} provide engineering frameworks for agent composition,
but not a formal semantics of global consistency.

Closer in spirit are attribution and coalition analyses based on Shapley values
or influence scores; these quantify contribution but not path-dependent trust.
Formal methods for distributed systems, proof-carrying code, and type systems
also contribute relevant ideas, but typically reason about program states rather
than epistemic transport between evidence-bearing judgments.  The contribution of
Judgment Geometry is precisely to combine the sheaf-theoretic language of local
and global compatibility with the differential-geometric language of connection,
curvature, and holonomy.

\section{Conclusion and future work}\label{sec:conclusion}

We have argued for a simple but powerful reframing: in multi-agent AI systems,
trust is not a score but a connection.  Once that move is made, a broad range of
otherwise slippery phenomena become computable.  Echo chambers become positive
curvature.  Hidden monopolies become holonomy.  Flat-but-ungrounded consensus
becomes a phantom region.  Existing JuGeo cohomology becomes the logical shadow
of a richer geometric structure.

The implementation work already shows that this perspective scales across the
JuGeo ecosystem.  The same ideas now appear in the core site verifier, the
multi-agent orchestration runtime, and the web verification stack.  In every
case, bundle geometry improves the system because it identifies structure that
local checks miss.

Several directions are immediately open.  Higher Chern classes could capture
more subtle topologies in large agent teams.  Spectral sequences could model
convergence over many rounds of interaction.  A Yang--Mills-like optimization
principle could be used to learn or synthesize optimal trust connections under
cost constraints.  And from an engineering perspective, bundle-aware routing and
repair policies may become one of the most practical ways to make AI agents both
more autonomous and more trustworthy.

\appendix

\section{Implementation notes}\label{app:impl}

The concrete implementation added during this work includes the following
artifacts.
\begin{itemize}
  \item \texttt{jugeo\_agents.core.bundle.JudgmentBundle}: computes trust
    connections, curvature, holonomy, the first Chern class, and trust
    stratification over agent outputs.
  \item \texttt{jugeo\_agents.core.fusion.GlobalSectionAssembler}: now carries
    bundle diagnostics alongside sheaf-theoretic fusion data.
  \item \texttt{jugeo\_agents.wrapper.JuGeoAgentWrapper}: now accumulates
    bundle diagnostics across the full pipeline and exposes them in
    \texttt{PipelineReport}.
  \item \texttt{jugeo.geometry.bundle.VerificationBundle}: brings the same
    construction to semantic-site verification in the JuGeo core library.
  \item \texttt{jugeo.webapp.trust.bundle.WebTrustBundle} plus
    \texttt{FiberedBundleDiagnostics} and \texttt{DeviceBundleDiagnostics}:
    extends connection geometry to request lifecycles, language fibers, and
    visual invariants across devices.
\end{itemize}

These modules collectively make bundle geometry a first-class concept across the
entire ecosystem rather than a standalone research note.

\begin{thebibliography}{99}

\bibitem{husemoller1994}
Dale Husemoller.
\newblock \emph{Fibre Bundles}.
\newblock Springer, third edition, 1994.

\bibitem{milnorstasheff1974}
John Milnor and James D. Stasheff.
\newblock \emph{Characteristic Classes}.
\newblock Princeton University Press, 1974.

\bibitem{botttu1982}
Raoul Bott and Loring W. Tu.
\newblock \emph{Differential Forms in Algebraic Topology}.
\newblock Springer, 1982.

\bibitem{goreskymacpherson1988}
Mark Goresky and Robert MacPherson.
\newblock \emph{Stratified Morse Theory}.
\newblock Springer, 1988.

\bibitem{sga4}
Michael Artin, Alexandre Grothendieck, and Jean-Louis Verdier.
\newblock \emph{Th\'eorie des topos et cohomologie \`etale des sch\'emas (SGA 4)}.
\newblock Lecture Notes in Mathematics, Springer, 1972--1973.

\bibitem{irving2018debate}
Geoffrey Irving, Paul Christiano, and Dario Amodei.
\newblock AI safety via debate.
\newblock \emph{arXiv preprint arXiv:1805.00899}, 2018.

\bibitem{wu2023autogen}
Qingyun Wu, Gagan Bansal, Jieyu Zhang, et al.
\newblock AutoGen: Enabling next-gen LLM applications via multi-agent conversation.
\newblock \emph{arXiv preprint arXiv:2308.08155}, 2023.

\bibitem{hong2023metagpt}
Sirui Hong, Xinlin Zheng, Jonathan Chen, et al.
\newblock MetaGPT: Meta programming for multi-agent collaborative framework.
\newblock \emph{arXiv preprint arXiv:2308.00352}, 2023.

\bibitem{crewai2024}
CrewAI.
\newblock CrewAI framework documentation.
\newblock \url{https://github.com/crewAIInc/crewAI}, accessed 2026.

\bibitem{langgraph2024}
LangChain.
\newblock LangGraph documentation.
\newblock \url{https://github.com/langchain-ai/langgraph}, accessed 2026.

\bibitem{jugeo-series}
JuGeo Research Group.
\newblock Judgment Geometry paper series.
\newblock Internal series including semantic sites, trust algebra, descent
obstructions, fleet competition, and agent verification, 2024--2026.

\bibitem{maclane1998}
Saunders Mac~Lane.
\newblock \emph{Categories for the Working Mathematician}.
\newblock Springer, second edition, 1998.

\bibitem{spivak2014}
David I. Spivak.
\newblock \emph{Category Theory for the Sciences}.
\newblock MIT Press, 2014.

\end{thebibliography}

\end{document}
